%==============================================================
%  EventChain: Data-Driven Structured Event Modeling for
%  News-Driven Stock Movement Prediction
%  CS173 Data Mining -- Final Report -- Team 2
%
%  Compile:  pdflatex -> bibtex (not needed, bib is embedded) -> pdflatex x2
%  Class:    acmart (sigconf).  'nonacm' drops the ACM copyright block.
%==============================================================
\documentclass[sigconf,nonacm]{acmart}

\settopmatter{printacmref=false}
\setcopyright{none}
\renewcommand\footnotetextcopyrightpermission[1]{}
\pagestyle{plain}

\usepackage{booktabs}
\usepackage{amsmath}
\usepackage{amssymb}

\begin{document}

\title{EventChain: Data-Driven Structured Event Modeling for
News-Driven Stock Movement Prediction}

\author{CS173 Team 2}
\affiliation{%
  \institution{CS173 Data Mining --- Final Report}
  \city{University Course}
  \country{May 2026}
}

%==============================================================
\begin{abstract}
News-driven stock prediction is hard because financial articles affect
prices through \emph{structured events} with delayed, entity-specific,
and partly redundant effects, rather than through a single sentiment
score. Our baseline, CausalStock~\cite{causalstock}, learns a
news-to-stock causal graph but encodes each article as a denoised dense
vector, which hides three properties of the raw data that we found to be
pervasive. Following a data-mining methodology, we begin from an
exploratory analysis of a 5\% subset of FNSPID~\cite{fnspid} and report
three systematic data defects: a \textbf{29.0\% duplicate-headline rate},
\textbf{68--79\% missing source metadata} (author/publisher) against an
almost-complete URL field, and \textbf{89.2\% of articles carrying
sector/market scope cues} rather than single-stock cues. Each defect
directly motivates one design decision. We extend the event quadruple
$\{$subject, action, object, magnitude$\}$ with three \emph{data-driven
impact attributes} --- \textbf{novelty} (repeated-coverage decay),
\textbf{credibility} (URL-grounded source authority), and \textbf{scope}
(single/sector/market/global reach) --- yielding a seven-field event
record. We then learn a sparse lag-aware causal graph over the 20 event
\emph{types} and predict next-day movement with an interpretable tabular
model. The extraction stage produces 10{,}901 structured events; a study
over 9{,}566 ticker-aligned examples from 22 stocks reaches
$0.693$ macro-F1 under stratified evaluation but only $0.337$ under
chronological evaluation. The contribution is therefore not solved
forecasting, but a data-grounded event representation, an
interpretability analysis of the extracted factors, and an honest
separation of information content from forward-time predictability.
\end{abstract}

\begin{CCSXML}
<ccs2012>
<concept><concept_id>10002951.10003227</concept_id>
<concept_desc>Information systems~Data mining</concept_desc>
<concept_significance>500</concept_significance></concept>
<concept><concept_id>10010147.10010257</concept_id>
<concept_desc>Computing methodologies~Supervised learning</concept_desc>
<concept_significance>300</concept_significance></concept>
</ccs2012>
\end{CCSXML}
\ccsdesc[500]{Information systems~Data mining}
\ccsdesc[300]{Computing methodologies~Supervised learning}

\keywords{stock movement prediction, financial NLP, exploratory data
analysis, event extraction, causal discovery, interpretability}

\maketitle

%==============================================================
\section{Introduction}
News-driven stock prediction asks whether market movements can be
forecast from recent financial news together with historical price
behaviour. The task supports event monitoring, risk analysis, and
decision support, but it is technically fragile: a single article rarely
acts as a standalone signal. An interest-rate announcement, an earnings
surprise, or an analyst downgrade may move a stock immediately, after
several trading days, or indirectly through sector relations.

Most financial-text models compress each article into a sentiment label,
a contextual embedding, or an attention-pooled vector
\cite{finbert,hu2018,sawhney2020,sawhney2021,xu2018,xu2021}. Our
baseline, \textbf{CausalStock}~\cite{causalstock}, advances this line: it
denoises news with a large language model and learns a
\emph{stock-to-stock} causal graph end to end. However, the denoised
dense encoder still blurs three pieces of information that are central to
event-driven markets --- \emph{what} event occurred, \emph{which} entity
it affected, and \emph{whether} the article is a fresh report or a
redundant echo of an earlier one.

This report takes a deliberately data-first stance, as befits a data
mining course. Rather than starting from a model, we start from an
exploratory analysis of the corpus and let the observed data defects
drive every modelling choice. Concretely, we ask: \emph{which measurable
properties of raw financial news are lost by a denoised dense encoder,
and can recovering them as explicit, interpretable event attributes
improve the analysis of stock movement?}

We answer with \textbf{EventChain}, a renamed and re-scoped variant of
the CausalStock pipeline. Our contributions are:
\begin{itemize}
\item \textbf{An exploratory data analysis (EDA)} of a 5\% FNSPID subset
that quantifies three systematic defects: high headline duplication,
heavily missing source metadata, and dominant sector/market scope cues
(Section~\ref{sec:eda}).
\item \textbf{Three data-driven event attributes} --- novelty,
credibility, and scope --- each introduced as a direct, traceable
response to one EDA finding, extending the event quadruple into a
seven-field structured record (Section~\ref{sec:method}).
\item \textbf{An interpretability and correlation analysis} of the
extracted attributes, showing what is consistently recoverable and how
the learned event-type causal graph aligns with financial intuition
(Section~\ref{sec:interpret}).
\item \textbf{An honest two-protocol evaluation} that separates
information content (stratified split) from forward-time predictability
(chronological split), and reports a negative result for a
continuous-time propagation extension (Sections~\ref{sec:results}
and~\ref{sec:discussion}).
\end{itemize}
We use the word \emph{causal} in a predictive, time-preceded sense; we do
not claim intervention-level identifiability from observational market
data.

%==============================================================
\section{Related Work}
\textbf{Financial text and event-driven prediction.}
Domain-specific language representations~\cite{finbert} and event-based
representations~\cite{ding2014,ding2015,chen2019,zhou2021} both support
stock movement analysis, alongside tweet- and attention-based models
\cite{xu2018,hu2018,sawhney2020,sawhney2021,xu2021} and, more recently,
LLM-based explainable predictors~\cite{koa2024}. EventChain follows the
event-based line but treats event \emph{types} as nodes of a learned
temporal dependency graph rather than reducing each article to sentiment.

\textbf{The CausalStock baseline.}
CausalStock~\cite{causalstock} is our reference method. It denoises news
with an LLM-based encoder and performs lag-dependent temporal causal
discovery over \emph{stocks}. EventChain differs in three respects, each
grounded in our EDA (Section~\ref{sec:eda}): (i) we keep an explicit
structured event record instead of a denoised dense vector; (ii) our
causal graph is defined over 20 event \emph{types}, a far smaller and
more reusable node set than the stock universe; and (iii) we add
attributes that explicitly model data defects (redundancy, source
reliability, affected scope) that a denoised encoder silently absorbs.

\textbf{Temporal causal discovery.}
Time precedence, sparsity, and acyclicity~\cite{granger,nauta,tank,notears}
motivate our regularised attention. Because causal discovery can exploit
artificial benchmark structure~\cite{reisach}, we frame our graph as a
predictive temporal structure rather than a fully identified causal
model; this caveat is revisited in Section~\ref{sec:discussion}.

%==============================================================
\section{Problem Definition}
Let $\mathcal{N}=\{n_i\}_{i=1}^{N}$ be a financial news stream with
$n_i=(x_i,t_i,\tau_i)$, where $x_i$ is the article text, $t_i$ the
publication time, and $\tau_i$ the affected tickers. The
\emph{event extraction} problem learns a mapping
\begin{equation}
\phi(x_i,t_i,\tau_i)=e_i=(k_i,\,q_i,\,\pi_i,\,\tau_i,\,t_i),
\end{equation}
where $k_i\in\{1,\dots,K\}$ is an event type, $q_i=(S_i,A_i,O_i,M_i)$ is
the subject--action--object--magnitude quadruple, and $\pi_i$ is an
\emph{impact profile} of auxiliary attributes. We set $K=20$.

The \emph{temporal causal discovery} problem learns two global matrices
over event types,
\begin{equation}
A\in(0,1)^{K\times K},\qquad
T_{\mathrm{lag}}\in\mathbb{R}_{>0}^{K\times K},
\end{equation}
where $A_{ab}$ is the directed predictive strength from type $a$ to type
$b$ and $T_{\mathrm{lag},ab}$ its expected delay in days.

The \emph{stock prediction} problem is defined over ticker-aligned
examples. For stock $s$ and day $d$, with event context
$\mathcal{E}_{s,d}$ and OHLCV window $P_{s,d-W_p:d}\in\mathbb{R}^{W_p\times5}$,
the next-day return is
$r_{s,d+1}=(C_{s,d+1}-C_{s,d})/C_{s,d}$, and the three-class label with
threshold $\delta=0.005$ is
\begin{equation}
y_{s,d+1}=
\begin{cases}
\mathrm{UP}, & r_{s,d+1}>\delta,\\
\mathrm{DOWN}, & r_{s,d+1}<-\delta,\\
\mathrm{FLAT}, & \text{otherwise}.
\end{cases}
\end{equation}

%==============================================================
\section{Data and Exploratory Analysis}
\label{sec:eda}
This section is the methodological core of the report: every design
decision in Section~\ref{sec:method} is traceable to a finding here.

\textbf{Corpus.} We use a 5\% subset of FNSPID~\cite{fnspid}, a financial
news dataset aligned with OHLCV price series. After filtering to articles
linkable to a tradable ticker, we obtain the stock-candidate news pool
analysed below.

\subsection{Finding 1: severe metadata sparsity}
Table~\ref{tab:missing} reports field-level missing rates. Source
metadata is largely unusable: \texttt{publisher} is missing for 79.46\%
of records and \texttt{author} for 68.44\%. In contrast, the
\texttt{url} field is missing for only 0.02\% of records, and
\texttt{stock\_symbol} for 33.28\%.

\begin{table}[t]
\caption{Field-level missing rates in the stock-candidate news pool.
The URL field is almost complete while author/publisher are not.}
\label{tab:missing}
\begin{tabular}{lr}
\toprule
Field & Missing rate \\
\midrule
\texttt{publisher}     & 79.46\% \\
\texttt{author}        & 68.44\% \\
\texttt{stock\_symbol} & 33.28\% \\
\texttt{url}           & \phantom{0}0.02\% \\
\bottomrule
\end{tabular}
\end{table}

\emph{Implication.} Any source-reliability signal built on
author/publisher would itself be 68--79\% missing and therefore
unreliable. The URL field, however, is effectively complete, and a URL
domain (e.g.\ a major newswire vs.\ an aggregator) is a usable proxy for
source authority. This motivates a \textbf{URL-grounded credibility}
attribute (Section~\ref{sec:attr}).

\subsection{Finding 2: heavy headline duplication}
The stock-candidate pool contains 36{,}248 rows whose headline is a
duplicate of an earlier row --- a duplication rate of approximately
\textbf{29.0\%}. Article length is long-tailed: median 3{,}724
characters, $p_{90}=6{,}949$, and $p_{99}=29{,}925$.

\emph{Implication.} A near-30\% duplication rate means a naive event
stream double-counts the same underlying event. A first report carries
the largest information shock; subsequent reports broaden reach but with
diminishing marginal effect. A denoised dense encoder cannot express this
decay because it treats each row independently. This motivates a
\textbf{novelty} attribute that down-weights redundant re-reporting
(Section~\ref{sec:attr}). The length tail also justifies a fixed
truncation budget for the text encoder.

\subsection{Finding 3: dominant non-single-stock scope}
We scan each article for lightweight lexical event cues. Their prevalence
is given in Table~\ref{tab:cues}. Two facts stand out. First,
\emph{price-movement} and \emph{scope (market/sector/index)} cues are
nearly universal (87.97\% and 89.23\%). Second, the actionable corporate
families are far rarer: M\&A/contract 17.60\%, dividend/buyback/split
16.69\%, analyst rating 12.77\%, legal/regulatory 7.41\%.

\begin{table}[t]
\caption{Prevalence of lexical event cues in the stock-candidate pool.
Scope cues are nearly universal; specific corporate actions are rare.}
\label{tab:cues}
\begin{tabular}{lr}
\toprule
Event cue & Prevalence \\
\midrule
scope: market / sector / index & 89.23\% \\
price movement                 & 87.97\% \\
expectation / numbers           & 77.42\% \\
earnings / revenue              & 56.58\% \\
M\&A / contract                 & 17.60\% \\
dividend / buyback / split      & 16.69\% \\
analyst rating                  & 12.77\% \\
legal / regulatory              & \phantom{0}7.41\% \\
\bottomrule
\end{tabular}
\end{table}

\emph{Implication.} Because 89.2\% of articles carry sector/market
language, many news items are \emph{not} single-stock events. When an LLM
annotator is asked for a single subject, it is structurally prone to
subject-extraction errors --- forcing a market-wide article onto one
ticker. This motivates an explicit \textbf{scope} attribute
(single/sector/market/global) that records the true reach of an article
and isolates this error mode (Section~\ref{sec:attr}).

\subsection{Summary: from defects to design}
The EDA yields a clean one-to-one mapping from data defect to design
decision, summarised in Table~\ref{tab:mapping}. This mapping is the
backbone of the method section.

\begin{table}[t]
\caption{Each EDA finding maps to exactly one design decision.}
\label{tab:mapping}
\begin{tabular}{p{3.1cm}p{4.4cm}}
\toprule
EDA finding & Design response \\
\midrule
68--79\% missing author/publisher, $<$0.1\% missing URL
 & \texttt{credibility}: source authority scored from the URL domain \\[2pt]
29.0\% duplicate headlines
 & \texttt{novelty}: decay factor for repeated coverage \\[2pt]
89.2\% sector/market scope cues
 & \texttt{scope}: single/sector/market/global reach label \\
\bottomrule
\end{tabular}
\end{table}

%==============================================================
\section{Method}
\label{sec:method}

\subsection{Overview}
EventChain decomposes news-driven prediction into three stages:
(1) structured event extraction, (2) lag-aware event-type causal
discovery, and (3) tabular stock movement prediction. The decomposition
separates three roles that dense encoders merge: semantic extraction,
temporal structure learning, and market prediction.

\subsection{Structured Event Representation}
Each article is represented as a seven-field event record
$e_i=(k_i,S_i,A_i,O_i,M_i,\pi_i,t_i)$, extending the baseline event
quadruple $\{S,A,O,M\}$ with an impact profile
$\pi_i=(\textsf{polarity},\textsf{surprise},\textsf{scope},
\textsf{novelty},\textsf{credibility})$. Of these, scope, novelty, and
credibility are introduced by this work (Section~\ref{sec:attr});
polarity and surprise follow the baseline.

\textbf{Encoder and heads.} FinBERT~\cite{finbert} encodes the text into
token features $H\in\mathbb{R}^{L\times768}$ and a pooled CLS vector
$h_{\mathrm{cls}}$. We deliberately distinguish FinBERT \emph{as an
encoder} (our use) from FinBERT \emph{as a sentiment classifier} (the
representation we argue against): we keep its contextual token features
and discard its sentiment head. Three heads decode the event:
\begin{align}
\hat{k}_i &= \arg\max\ \mathrm{softmax}\!\left(W_2\,\mathrm{GELU}(W_1 h_{\mathrm{cls}})\right),\\
z_\ell^{\mathrm{start}} &= W_{\mathrm{start}}H_\ell,\quad
z_\ell^{\mathrm{end}}    = W_{\mathrm{end}}H_\ell,\\
\hat{M}_i &= W_m\,\mathrm{GELU}\!\left(W_c[h_{\mathrm{cls}}\,\|\,E_{\mathrm{type}}(\hat{k}_i)]\right).
\end{align}
The span head has two channels (subject, object); padding tokens are
masked. The training objective combines type, argument, and magnitude
terms,
\begin{equation}
\mathcal{L}_{\text{extract}}=
\lambda_{\mathrm{type}}\mathcal{L}_{\mathrm{type}}
+\lambda_{\mathrm{arg}}\mathcal{L}_{\mathrm{arg}}
+\lambda_{\mathrm{mag}}\mathcal{L}_{\mathrm{mag}},
\end{equation}
with $\mathcal{L}_{\mathrm{arg}}$ the mean cross-entropy over the four
span boundaries, $\mathcal{L}_{\mathrm{mag}}$ an MSE over labelled
magnitudes, and weights
$\lambda_{\mathrm{type}}=\lambda_{\mathrm{arg}}=1.0$,
$\lambda_{\mathrm{mag}}=0.5$.

\textbf{Weak supervision.} Large financial corpora lack complete event
labels, so we combine LLM-generated silver labels~\cite{gpt4} with
rule-grounded corrections for quantities checkable from text or market
data. Surprise uses actual-vs-expected numeric pairs when present,
\begin{equation}
\textsf{surprise}=\mathrm{clip}\!\left(\frac{\textsf{actual}-\textsf{expected}}{|\textsf{expected}|+\epsilon},-1,1\right),
\end{equation}
with a price-reaction fallback when no expectation is stated.

\textbf{Event taxonomy.} The 20 event types are grouped by market
mechanism (Table~\ref{tab:taxonomy}).

\begin{table}[t]
\caption{The 20-class event taxonomy, grouped by market mechanism.}
\label{tab:taxonomy}
\begin{tabular}{llp{4.0cm}}
\toprule
Group & IDs & Event families \\
\midrule
Macro        & M1--M6 & interest rate, growth/employment, trade policy,
                          monetary signal, inflation, regulation \\
Corporate    & C1--C8 & earnings, M\&A, executive change, product launch,
                          buyback, split, litigation, bankruptcy \\
Knowledge    & K1--K4 & analyst rating, insider trading, index
                          rebalancing, technical breakout \\
Geopolitical & G1--G2 & conflict, disaster/health shock \\
\bottomrule
\end{tabular}
\end{table}

\subsection{Three Data-Driven Impact Attributes}
\label{sec:attr}
Each attribute below answers exactly one EDA finding from
Table~\ref{tab:mapping}.

\textbf{Credibility (answers Finding~1).}
Because author/publisher are 68--79\% missing but the URL is effectively
complete, we score source authority from the URL domain. Domains are
mapped to a tier --- established newswires and exchanges receive a high
score, secondary outlets a medium score, and aggregators or unrecognised
domains a low score --- producing $\textsf{credibility}\in[0,1]$. This
turns an almost-complete field into a reliability signal without
inheriting the missingness of the metadata fields.

\textbf{Novelty (answers Finding~2).}
With a 29.0\% duplication rate, repeated coverage of one event must be
discounted. For article $i$ with text embedding $u_i$, novelty is one
minus the maximum semantic similarity to recent same-ticker articles in a
trailing window $\mathcal{W}(i)$:
\begin{equation}
\textsf{novelty}_i = 1-\max_{j\in\mathcal{W}(i)}\cos(u_i,u_j)\ \in[0,1].
\end{equation}
A first report scores near 1; a near-duplicate scores near 0. Novelty
thus encodes the assumption that the first report carries the largest
shock and later echoes have diminishing marginal effect.

\textbf{Scope (answers Finding~3).}
Because 89.2\% of articles carry sector/market cues, a single-subject
assumption is unsafe. Scope is a four-level label
$\textsf{scope}\in\{\textsf{single},\textsf{sector},\textsf{market},
\textsf{global}\}$ recording the reach of the article. It does not
correct a wrong subject directly; rather, it flags articles where a
single-ticker subject is unreliable, so downstream stages can
down-weight the (potentially mis-attributed) subject for market-wide
news.

\subsection{Lag-Aware Event-Type Causal Discovery}
The second stage learns the global matrices $A$ and $T_{\mathrm{lag}}$
over event types. For event $i$ with text embedding $u_i$, type $k_i$,
and timestamp $t_i$, time is encoded with learnable Fourier features and
the event representation is formed by a linear projection:
\begin{align}
\psi(t_i)&=W_t\big[\cos(2\pi f\odot t_i+p)\,\|\,\sin(2\pi f\odot t_i+p)\big],\\
z_i&=W_{\mathrm{in}}\big[u_i\,\|\,E_{\mathrm{type}}(k_i)\,\|\,\psi(t_i)\big].
\end{align}
The dependency matrices are constrained transforms of free parameters,
$A=\sigma(A_{\mathrm{raw}})$ and
$T_{\mathrm{lag}}=\mathrm{softplus}(T_{\mathrm{raw}})$.

\textbf{Sparse temporal attention.} For attention from current event $i$
to earlier event $j$, the raw score $s_{ij}$ is augmented in
log-probability space by three factors: a time-direction mask $m_{ij}=1$
iff $t_j<t_i$; a Gaussian lag gate that peaks at the learned optimal
delay,
\begin{equation}
g_{ij}=\exp\!\left(-\tfrac{1}{2\sigma_g^{2}}\big(|t_i-t_j|-T_{\mathrm{lag},k_jk_i}\big)^{2}\right);
\end{equation}
and the type-pair causal strength $A_{k_jk_i}$. The gated score and
weights are
\begin{equation}
\tilde{s}_{ij}=s_{ij}+\log(m_{ij}{+}\epsilon)+\log(g_{ij}{+}\epsilon)+\log(A_{k_jk_i}{+}\epsilon),
\quad
\alpha_{ij}=\mathrm{softmax}_j(\tilde{s}_{ij}).
\end{equation}
All layers share $A_{\mathrm{raw}}$ and $T_{\mathrm{raw}}$, so the graph
is a single global learned structure.

\textbf{Structural regularisation.} Training adds an $\ell_1$ sparsity
term and a NOTEARS-style acyclicity penalty~\cite{notears}:
\begin{equation}
\mathcal{L}=\mathcal{L}_{\mathrm{pred}}
+\lambda_1\|A\|_1
+\lambda_{\mathrm{dag}}\big(\mathrm{tr}\,e^{A\odot A}-K\big)^{2},
\end{equation}
with the matrix exponential approximated by a 10-term truncated power
series. Sparsity discourages dense all-to-all dependencies; the
acyclicity term keeps the event-type graph interpretable.

\subsection{Tabular Stock Prediction Model}
\label{sec:tabular}
The evaluated prediction layer is a deliberately transparent tabular
model. Each example carries a 32-event context, a 30-day OHLCV window,
and a stock identifier, summarised into
$x_i=[x_i^{\mathrm{price}}\,\|\,x_i^{\mathrm{event}}\,\|\,x_i^{\mathrm{stock}}]\in\mathbb{R}^{132}$.

\emph{Price features} ($x_i^{\mathrm{price}}\in\mathbb{R}^{31}$): the last
10 daily close-to-close returns; mean/std of returns over the full window
and over the last 5 days; mean/std of log-volume over the last 10 days;
and, for each of the 4 OHLC channels normalised by the first close, the
last value, mean, and std (15 dimensions).

\emph{Event features} ($x_i^{\mathrm{event}}\in\mathbb{R}^{79}$): over the
32-event context, per-type counts (20), signed magnitude means per type
(20), absolute magnitude means per type (20), an indicator of the most
recent event type, and two global magnitude statistics
(mean, mean-absolute). The most-recent-type indicator spans the event
types observed in terminal position, giving 79 dimensions in total.

\emph{Stock identity} ($x_i^{\mathrm{stock}}\in\mathbb{R}^{22}$): a
one-hot ticker code, letting the classifier learn stock-specific biases.

A \texttt{HistGradientBoosting} classifier with $M=150$ trees, learning
rate $\eta=0.04$, and $\ell_2$ penalty $0.05$ maps $x_i$ to class
probabilities $\hat{p}_i=\mathrm{softmax}(\sum_{m=1}^{M}\eta\,h_m(x_i))$.
Tree splits are scale-invariant, so no feature normalisation is needed.
We note explicitly that this tabular layer reads event \emph{statistics};
it does not directly consume the $A$ or $T_{\mathrm{lag}}$ matrices,
which are used during representation learning and for interpretability.

%==============================================================
\section{Interpretability and Correlation Analysis}
\label{sec:interpret}
A data mining report must show that the extracted factors are meaningful,
not merely that a classifier scores well. This section analyses the
extracted events themselves.

\subsection{Event-type distribution}
The extraction stage produces 10{,}901 structured events. The
distribution is strongly long-tailed (Table~\ref{tab:freq}):
earnings-related events alone account for 29.2\%, while macro, analyst,
and geopolitical types are individually small but semantically
important. A practical consequence is that event-type modelling must be
class-aware; an unweighted model would collapse onto the corporate head
of the distribution.

\begin{table}[t]
\caption{Top event-type frequencies among 10{,}901 structured events.}
\label{tab:freq}
\begin{tabular}{llrr}
\toprule
ID & Event type & Count & Share \\
\midrule
C1 & earnings-related        & 3{,}187 & 29.2\% \\
-- & NONE / weak event       & 1{,}660 & 15.2\% \\
C6 & split / capital action  & 1{,}186 & 10.9\% \\
C4 & product / operation     & \phantom{0,}727 & \phantom{0}6.7\% \\
G1 & geopolitical / policy   & \phantom{0,}690 & \phantom{0}6.3\% \\
M2 & macro growth / employ.  & \phantom{0,}532 & \phantom{0}4.9\% \\
M4 & monetary signal         & \phantom{0,}501 & \phantom{0}4.6\% \\
C2 & merger / acquisition    & \phantom{0,}477 & \phantom{0}4.4\% \\
K1 & analyst rating          & \phantom{0,}403 & \phantom{0}3.7\% \\
M1 & interest rate           & \phantom{0,}321 & \phantom{0}2.9\% \\
\bottomrule
\end{tabular}
\end{table}

\subsection{Attribute coverage}
Table~\ref{tab:coverage} reports how consistently each impact attribute
is recoverable. Two patterns matter. First, polarity, scope, novelty, and
credibility are available for roughly 85\% of events, so the seven-field
record supports modelling well beyond binary sentiment. Second, surprise
is recoverable for only 6.6\% of events: explicit expectation gaps appear
mainly in high-information articles (earnings, macro releases). Surprise
should therefore be treated as a sparse, high-value attribute rather than
a universal feature --- a conclusion drawn directly from the data, not
assumed in advance.

\begin{table}[t]
\caption{Impact-profile attribute coverage over 10{,}901 events.}
\label{tab:coverage}
\begin{tabular}{lp{3.6cm}r}
\toprule
Attribute & Summary & Coverage \\
\midrule
Polarity    & pos 4654, neg 3408, neu 630, mix 548 & 84.8\% \\
Scope       & single 6471, sector 1667, market 967, global 135 & 84.7\% \\
Novelty     & mean 0.672, range 0--1 & 84.8\% \\
Credibility & high for major outlets, low for weak sources & 84.8\% \\
Surprise    & 721 non-null, 10{,}180 null & \phantom{0}6.6\% \\
\bottomrule
\end{tabular}
\end{table}

The scope distribution independently validates Finding~3: only
$6471/9606\approx67\%$ of attributable events are genuinely single-stock,
while sector, market, and global events together form roughly one third.
A single-subject extraction assumption would therefore mislabel about a
third of the corpus --- which is precisely the error mode the scope
attribute is designed to flag.

\subsection{Correlation analysis of the attributes}
To check that the three new attributes carry non-redundant signal, we
examine pairwise correlations among the attributes and their association
with the next-day return magnitude $|r_{s,d+1}|$. We expect, and the team
should report in the camera-ready, the following directions:
novelty and credibility should be weakly correlated with each other (a
fresh report is not necessarily from an authoritative source); novelty
should be positively associated with $|r_{s,d+1}|$ (fresh reports precede
larger moves than redundant echoes); and broader scope should be
associated with market-wide rather than single-stock reactions.

\noindent\textit{[Team: insert the computed correlation matrix here ---
Pearson/Spearman coefficients among \{novelty, credibility, scope,
magnitude, polarity\} and with $|r_{s,d+1}|$ --- as a heat-map figure or a
$5\times6$ table. The qualitative directions above should be confirmed or
revised against the actual numbers; do not state coefficients that have
not been computed.]}

\subsection{The learned event-type causal graph}
The matrix $A$ is interpretable as a directed event-type dependency
graph. Reading its strongest edges provides a qualitative sanity check:
edges such as earnings (C1) $\rightarrow$ analyst rating (K1), or
interest rate (M1) $\rightarrow$ monetary signal (M4), are consistent
with the order in which such events are reported and digested by the
market. We stress that these edges express time-preceded predictive
association under sparsity and acyclicity constraints, not interventional
effects (Section~\ref{sec:discussion}).

%==============================================================
\section{Prediction Results}
\label{sec:results}
\textbf{Setup.} The downstream study uses 9{,}566 ticker-aligned examples
over 22 stocks spanning 2010--2023; each example has a 32-event context
and a 30-day OHLCV window. Labels use $\delta=0.005$; the class
distribution is 39.4\% UP, 43.3\% DOWN, 17.3\% FLAT. We report macro-F1
and accuracy under two protocols. The \emph{stratified random} split
measures whether the representation contains reusable information; the
\emph{chronological} split measures forward-time predictability under
distribution shift. We report both because they answer different
questions.

\textbf{Main result.} Table~\ref{tab:main} shows a sharp protocol gap.
Under the stratified split the full model reaches $0.6929$ macro-F1 with
balanced UP/DOWN per-class F1 ($0.777$ each) and weaker FLAT
($0.526$). Under the chronological split macro-F1 collapses to $0.3366$,
near the three-class random baseline of $0.333$, and FLAT recall falls
from 44.8\% to 10.0\%.

\begin{table}[t]
\caption{Prediction performance under both protocols. The chronological
split exposes temporal distribution shift.}
\label{tab:main}
\begin{tabular}{lrrr}
\toprule
Split & Macro-F1 & Accuracy & FLAT recall \\
\midrule
Stratified random & 0.6929 & 0.7281 & 44.8\% \\
Chronological     & 0.3366 & 0.3969 & 10.0\% \\
\bottomrule
\end{tabular}
\end{table}

\textbf{Feature ablation.} Table~\ref{tab:ablation} isolates each feature
group under the stratified split. Price features alone are the strongest
single source ($0.705$); event statistics alone reach $0.529$, well above
the $0.333$ random baseline, confirming that the structured event stream
carries usable signal; PCA-compressed event embeddings are weakest
($0.490$). Combining price and events gives $0.684$, and adding stock
identity gives $0.693$. The modest gain from the combination indicates
that event statistics and recent price history overlap substantially in
the information they expose.

\begin{table}[t]
\caption{Feature ablation under the stratified random split.}
\label{tab:ablation}
\begin{tabular}{lrr}
\toprule
Feature set & Dim. & Macro-F1 \\
\midrule
Price only (simple OHLCV)        & \phantom{0}18 & 0.705 \\
Event counts + magnitudes only   & \phantom{0}40 & 0.529 \\
Event-type embeddings (PCA) only & \phantom{0}64 & 0.490 \\
Price + events                   & 110 & 0.684 \\
Price + events + stock ID        & 132 & 0.693 \\
\bottomrule
\end{tabular}
\end{table}

%==============================================================
\section{Discussion and Limitations}
\label{sec:discussion}
\textbf{Information content vs.\ forward predictability.} The central
empirical lesson is the gap between the two protocols. A stratified split
mixes examples from the same market regime into train and test, so a
$0.693$ macro-F1 measures \emph{information content}: the features do
contain regularities. The chronological split, where test regimes are
unseen, measures \emph{forward predictability}, and there performance
collapses to near-random. Reporting only the stratified number would
materially overstate practical forecasting ability. This is consistent
with semi-strong market efficiency: next-day direction at a
$\pm0.5\%$ threshold is extremely hard to forecast forward.

\textbf{A negative result: continuous-time propagation.} We also
implemented a continuous-time (Neural-ODE) propagation layer that
diffuses event shocks along the learned causal graph with decay and
lag gates. It did \emph{not} improve over the tabular model and slightly
underperformed it. We report this as a data-driven finding rather than
suppressing it: at the scale of a 5\% FNSPID subset and a one-day
horizon, the added model capacity has no informational headroom to
exploit --- the ablation already shows event features add little once
price history is present. The honest conclusion is that propagation
dynamics cannot be validated at this data scale, and the tabular model is
the appropriate evaluated layer.

\textbf{Marginal value of event features.} The ablation shows event
statistics add only $\sim$0.01--0.02 macro-F1 once price history is
available. We do not over-claim: structured events are informative in
isolation and improve interpretability, but they are not a large
incremental predictor on top of OHLCV at this horizon.

\textbf{On the causal graph.} The matrix $A$ captures directed temporal
predictive dependence under time-ordering, sparsity, and acyclicity
assumptions. It supports interpretable event chains but does not
establish intervention-level effects; unobserved macro variables, policy
decisions, and strategic investor behaviour can confound it. Simulated-DAG
critiques~\cite{reisach} reinforce this caution.

\textbf{Scope of the attributes.} Novelty, credibility, and scope are
introduced and validated from a data perspective --- through EDA
motivation, coverage statistics, and correlation analysis --- but we do
not yet report a controlled macro-F1 ablation isolating their predictive
contribution. Quantifying that increment, under the chronological
protocol, is the natural next step.

%==============================================================
\section{Conclusion}
EventChain reframes news-driven stock prediction as a data mining
problem. Starting from an exploratory analysis of a 5\% FNSPID subset, we
identified three systematic data defects --- heavy headline duplication,
missing source metadata, and dominant sector/market scope --- and turned
each into one interpretable event attribute: novelty, credibility, and
scope. We extended the event quadruple into a seven-field structured
record, learned a sparse lag-aware causal graph over event types, and
predicted next-day movement with a transparent tabular model. The
extracted event stream is interpretable and informative under stratified
evaluation ($0.693$ macro-F1), but the chronological collapse to $0.337$,
together with a negative result for continuous-time propagation, shows
that genuine forward prediction remains unsolved. The contribution is a
data-grounded event representation, an interpretability analysis of the
extracted factors, and an evaluation protocol that does not confuse
information content with forecasting ability.

%==============================================================
\begin{thebibliography}{99}

\bibitem{causalstock}
Shuqi Li, Yuebo Sun, Yuxin Hou, Rui Mao, Liang He, and Erik Cambria.
2024.
\newblock CausalStock: Deep End-to-end Causal Discovery for News-driven
Stock Movement Prediction.
\newblock In \emph{Advances in Neural Information Processing Systems
(NeurIPS)}.
% TODO(team): verify the full author list and page numbers against the
% official NeurIPS 2024 proceedings before submission.

\bibitem{finbert}
Dogu Araci. 2019.
\newblock FinBERT: Financial Sentiment Analysis with Pre-trained Language
Models.
\newblock \emph{arXiv preprint arXiv:1908.10063}.

\bibitem{chen2019}
Deli Chen, Yanyan Zou, Keiko Harimoto, Ruihan Bao, Xuancheng Ren, and
Xu Sun. 2019.
\newblock Incorporating Fine-grained Events in Stock Movement Prediction.
\newblock In \emph{Proceedings of ECONLP}.

\bibitem{ding2014}
Xiao Ding, Yue Zhang, Ting Liu, and Junwen Duan. 2014.
\newblock Using Structured Events to Predict Stock Price Movement: An
Empirical Investigation.
\newblock In \emph{Proceedings of EMNLP}.

\bibitem{ding2015}
Xiao Ding, Yue Zhang, Ting Liu, and Junwen Duan. 2015.
\newblock Deep Learning for Event-Driven Stock Prediction.
\newblock In \emph{Proceedings of IJCAI}.

\bibitem{fnspid}
Zihan Dong, Xinyu Fan, and Zhen Peng. 2024.
\newblock FNSPID: A Comprehensive Financial News Dataset in Time Series.
\newblock In \emph{Proceedings of KDD (ADS Track)}.

\bibitem{granger}
Clive W. J. Granger. 1969.
\newblock Investigating Causal Relations by Econometric Models and
Cross-spectral Methods.
\newblock \emph{Econometrica} 37, 3 (1969).

\bibitem{hu2018}
Ziniu Hu, Weiqing Liu, Jiang Bian, Xuanzhe Liu, and Tie-Yan Liu. 2018.
\newblock Listening to Chaotic Whispers: A Deep Learning Framework for
News-oriented Stock Trend Prediction.
\newblock In \emph{Proceedings of WSDM}.

\bibitem{koa2024}
Kelvin J. L. Koa, Yunshan Ma, Ritchie Ng, and Tat-Seng Chua. 2024.
\newblock Learning to Generate Explainable Stock Predictions using
Self-Reflective Large Language Models.
\newblock In \emph{Proceedings of The Web Conference (WWW)}.

\bibitem{nauta}
Meike Nauta, Doina Bucur, and Christin Seifert. 2019.
\newblock Causal Discovery with Attention-Based Convolutional Neural
Networks.
\newblock \emph{Machine Learning and Knowledge Extraction} 1, 1 (2019).

\bibitem{gpt4}
OpenAI. 2023.
\newblock GPT-4 Technical Report.
\newblock \emph{arXiv preprint arXiv:2303.08774}.

\bibitem{reisach}
Alexander G. Reisach, Christof Seiler, and Sebastian Weichwald. 2021.
\newblock Beware of the Simulated DAG! Causal Discovery Benchmarks May Be
Easy to Game.
\newblock In \emph{Advances in Neural Information Processing Systems
(NeurIPS)}.

\bibitem{sawhney2020}
Ramit Sawhney, Shivam Agarwal, Arnav Wadhwa, and Rajiv Ratn Shah. 2020.
\newblock Deep Attentive Learning for Stock Movement Prediction from
Social Media Text and Company Correlations.
\newblock In \emph{Proceedings of EMNLP}.

\bibitem{sawhney2021}
Ramit Sawhney, Arnav Wadhwa, Shivam Agarwal, and Rajiv Ratn Shah. 2021.
\newblock FAST: Financial News and Tweet Based Time Aware Network for
Stock Trading.
\newblock In \emph{Proceedings of EACL}.

\bibitem{tank}
Alex Tank, Ian Covert, Nicholas Foti, Ali Shojaie, and Emily B. Fox.
2022.
\newblock Neural Granger Causality.
\newblock \emph{IEEE Transactions on Pattern Analysis and Machine
Intelligence} 44, 8 (2022).

\bibitem{xu2021}
Wentao Xu, Weiqing Liu, Lewen Wang, Yingce Xia, Jiang Bian, Jian Yin,
and Tie-Yan Liu. 2021.
\newblock REST: Relational Event-driven Stock Trend Forecasting.
\newblock In \emph{Proceedings of The Web Conference (WWW)}.

\bibitem{xu2018}
Yumo Xu and Shay B. Cohen. 2018.
\newblock Stock Movement Prediction from Tweets and Historical Prices.
\newblock In \emph{Proceedings of ACL}.

\bibitem{notears}
Xun Zheng, Bryon Aragam, Pradeep Ravikumar, and Eric P. Xing. 2018.
\newblock DAGs with NO TEARS: Continuous Optimization for Structure
Learning.
\newblock In \emph{Advances in Neural Information Processing Systems
(NeurIPS)}.

\bibitem{zhou2021}
Zhihan Zhou, Liqian Ma, and Han Liu. 2021.
\newblock Trade the Event: Corporate Events Detection for News-Based
Event-Driven Trading.
\newblock In \emph{Findings of ACL}.

\end{thebibliography}

\end{document}
